\subsection{A concordance obstruction over cyclic three-manifolds}

Let $B$ be a connected, closed, oriented real PL three-manifold.
Write $P_{\mathrm{TOP}}(B)$ for the compact-open space of
homeomorphisms of $B\times I$ that fix $B\times\{0\}$ pointwise;
the top face is free.  Write $P_{\mathrm{PL}}(B)$ for the corresponding
simplicial group of jointly PL families.  Thus a parameterized PL
family is required to be PL on its total space.

\begin{theorem}\label{thm:cyclic-three-concordance-obstruction}
Suppose that $\pi_1(B)$ is finite cyclic of even order.  Then there
is a homeomorphism
\[
 F:B\times I\longrightarrow B\times I,
 \qquad F|_{B\times\partial I}=\operatorname{id},
\]
whose component in $P_{\mathrm{TOP}}(B)$ contains no PL
concordance.  In particular,
\[
 \pi_0P_{\mathrm{PL}}(B)\longrightarrow
 \pi_0P_{\mathrm{TOP}}(B)
\]
is not surjective.  The same conclusion holds for the map from
definable concordance components, for every fixed o-minimal
expansion of $\mathbb R$.  The standard PL manifold
$\mathbb {RP}^{3}$ is an example.
\end{theorem}

We first describe the relative obstruction and the two classical
inputs used in the proof.  All cohomology groups without an
explicit coefficient group in this subsection have coefficients
$\mathbb Z/2$.

\paragraph{The formal PL difference class.}
If $X$ is a compact PL four-manifold and $f:X\to X$ fixes
$\partial X$ pointwise, its relative Casson--Sullivan class is
\begin{equation}\label{eq:relative-three-base-cs}
 \operatorname{cs}_{\mathrm{PL}}(f)
 =\sigma^{-1}\operatorname{ks}(M_f,\partial M_f)
 \in H^3(X,\partial X),\qquad
 M_f=(X\times I)\cup_f X.
\end{equation}
Here the two ends have the given PL tangent reductions, the
vertical boundary has their product reduction, and
$\operatorname{ks}$ is the obstruction to extending this
prescribed stable PL reduction across the five-dimensional
mapping cylinder.  The isomorphism
\[
 \sigma:H^3(X,\partial X)\xrightarrow{\cong}
             H^4(M_f,\partial M_f)
\]
is the relative suspension isomorphism.  Tangent microbundles
are identified using the canonical isomorphism induced by the
homeomorphism $f$.

This class measures the difference between the specified formal
stable tangent reductions.
Choose a smoothing compatible with the given PL structure,
whose existence is supplied by
\cite[Essay~I, p.~39]{KirbySiebenmann1977}.
The map $BO\to BPL$ identifies the degree-four
Kirby--Siebenmann obstruction, so \eqref{eq:relative-three-base-cs}
agrees with the smooth formal Casson--Sullivan class in
\cite[Remark~2.5 and Definition~2.11, pp.~2502--2504]{Galvin2026CassonSullivan}.
In particular it is invariant under an ordinary isotopy fixing
$\partial X$, and it vanishes for a PL homeomorphism fixing
$\partial X$.  The first assertion follows either by the
relative tangent-microbundle homotopy or by
\cite[Proposition~2.22, pp.~2507--2508]{Galvin2026CassonSullivan}.
For the second, the PL mapping cylinder extends the prescribed
boundary PL reduction over its entire interior, so its relative
obstruction vanishes.

\begin{lemma}\label{lem:compact-cyclic-relative-cs-realization}
Let $X$ be a compact, connected, oriented smooth four-manifold,
possibly with boundary, with finite cyclic fundamental group.
Every element of $H^3(X,\partial X)$ is
$\operatorname{cs}_{\mathrm{PL}}(f)$ for a homeomorphism
$f:X\to X$ which fixes $\partial X$ pointwise.
\end{lemma}

\begin{proof}
We spell out the boundary conditions in the normal-invariant
construction of
\cite[Section~4.2 and proof of Theorem~1.7, pp.~2524--2528]{Galvin2026CassonSullivan}.
Put $M=X\times I$ and retain the full boundary
\[
 \partial M=(X\times\partial I)\cup(\partial X\times I).
\]
Compatible collars, or a fixed rounding of corners, will be
understood.  The relative normal invariants have coordinates
\begin{equation}\label{eq:cyclic-relative-normal-coordinates}
 [M/\partial M,G/\mathrm{TOP}]
 \cong H^2(M,\partial M;\mathbb Z/2)
                  \oplus H^4(M,\partial M;\mathbb Z).
\end{equation}
Under the map to the relative stable smoothing obstruction,
the element $(0,y)$ has image $\operatorname{red}_2(y)$;
these are the coordinates and the characteristic-class formula
of Lemmas~4.8--4.9 in that source.  Only the zero first
coordinate is used here.

The coefficient exact sequence contains
\[
 H^4(M,\partial M;\mathbb Z)\xrightarrow{\operatorname{red}_2}
 H^4(M,\partial M;\mathbb Z/2)\longrightarrow
 H^5(M,\partial M;\mathbb Z)\xrightarrow{\,2\,}
 H^5(M,\partial M;\mathbb Z).
\]
The last integral group is $\mathbb Z$, by orientation and
Poincar\'e--Lefschetz duality.  Consequently reduction is
surjective.  Given $\eta\in H^3(X,\partial X)$, choose $y$
whose reduction is its suspension.

For a finite cyclic group $\pi$, Bak's calculation gives
$L_5^s(\mathbb Z[\pi])=0$, as recalled in
\cite[Section~4.4, p.~2528]{Galvin2026CassonSullivan}.
The surgery obstruction of $(0,y)$ therefore vanishes.
Exactness of the simple relative TOP surgery sequence
\cite[Theorem~11.3A and the simple-structure convention, p.~200]{FreedmanQuinn1990}
provides a simple homotopy equivalence of triads
\[
 \phi:(N;\partial_+N,\partial_0N,\partial_-N)
 \longrightarrow
 (X\times I;X\times\{0\},\partial X\times I,X\times\{1\}),
\]
whose restriction to the full boundary is a homeomorphism.
The smooth boundary markings are those induced by $\phi$.

Finite groups are good in the sense of Freedman--Quinn
\cite[p.~99]{FreedmanQuinn1990}.  The compact relative
five-dimensional $s$-cobordism theorem
\cite[Theorem~7.1A, p.~101]{FreedmanQuinn1990}, in the explicitly
marked form of
\cite[Proposition~4.10, pp.~2526--2527]{Galvin2026CassonSullivan},
gives a product homeomorphism
\[
 p:N\xrightarrow{\cong}\partial_+N\times I
\]
which is the identity on $\partial_+N$ and on its full vertical
boundary.  Write $p_-$ for its top restriction and $\phi_\pm$
for the end markings, regarded as maps to $X$.  Then
\[
 f=\phi_+\,p_-\,\phi_-^{-1}:X\longrightarrow X
\]
is the identity on $\partial X$, because the product structure
and the two end markings agree on the marked vertical boundary.
Moreover, $N$ is the mapping cylinder $M_f$ with exactly its
prescribed boundary reductions.  By the relative obstruction
identification in
\cite[Lemma~4.12, p.~2527]{Galvin2026CassonSullivan},
\[
 \operatorname{cs}_{\mathrm{PL}}(f)
   =\sigma^{-1}\operatorname{red}_2(y)=\eta.
\]
This proves the assertion with every boundary face fixed.
\end{proof}

\paragraph{Replacing an endpoint path by a jointly PL path.}
For a closed real PL three-manifold the automorphism comparison
\eqref{eq:three-manifold-full-automorphism-comparison} is a weak
equivalence.  Its parameter convention is the total-space PL
convention of
\cite[Section~2.1, p.~222; Theorem~3.1.2, pp.~225--226]{Hamstrom1974}.
We need only injectivity on components and surjectivity on
the fundamental group.

It follows that if $a_t$, $0\leq t\leq1$, is an ordinary path
of homeomorphisms with $a_0=\operatorname{id}$ and $a_1$ PL,
then there is a jointly PL path $b_t$ with these endpoints
which is ordinary homotopic to $a_t$ relative to its endpoints.
Indeed, component injectivity first supplies a jointly PL path
$b^0$ with the same endpoints.  The loop $a*(b^0)^{-1}$ is
ordinary homotopic to a jointly PL loop $\ell$, by surjectivity
on $\pi_1$.  Concatenating $\ell$ with $b^0$ gives the required
$b$.  Concatenations are made with piecewise affine
reparameterizations, which preserve the joint PL condition.

\begin{proof}[Proof of Theorem~\ref{thm:cyclic-three-concordance-obstruction}]
Give $B$ a smoothing compatible with its PL structure and put
$X=B\times I$.  Mod-two duality and the abelianization of its
fundamental group give
\[
 H^3(X,\partial X)\cong H^2(B)
 \cong H_1(B;\mathbb Z/2)\cong\mathbb Z/2.
\]
Apply Lemma~\ref{lem:compact-cyclic-relative-cs-realization} to
the nonzero class.  It gives a homeomorphism $F$ fixing both
ends with $\operatorname{cs}_{\mathrm{PL}}(F)\ne0$.

Suppose that an ordinary path $H_t$ in $P_{\mathrm{TOP}}(B)$
joins $F$ to a PL concordance $G$.  Its top restrictions
$a_t=H_t|_{B\times\{1\}}$ run from the identity to the PL
map $G|_{B\times\{1\}}$.  Define
\[
 T_t(x,s)=(a_{ts}(x),s),\qquad Q_t=T_t^{-1}H_t.
\]
These are ordinary TOP families; compactness of $B$ ensures
continuity of their inverse families.  Each $Q_t$ fixes both
ends.  Thus $Q_t$ is a relative-boundary isotopy from $F$ to
$T_1^{-1}G$.

Replace $a$ by the jointly PL path $b$ constructed above, and
write $T^b(x,s)=(b_s(x),s)$.  An endpoint-fixed homotopy from
$a$ to $b$ gives an isotopy of their traces with fixed boundary
restrictions.  After inversion and composition with $G$, it
gives an isotopy relative to the entire boundary from
$T_1^{-1}G$ to $(T^b)^{-1}G$.  The latter is a PL
homeomorphism fixing both ends.  Relative isotopy invariance
and PL vanishing in \eqref{eq:relative-three-base-cs} imply
$\operatorname{cs}_{\mathrm{PL}}(F)=0$, a contradiction.

For a fixed o-minimal expansion of $\mathbb R$, the relative
definable--PL component comparison of Theorem~\ref{result-4}
replaces every definable concordance by a PL one through a
definable isotopy with the bottom fixed pointwise.  Its
underlying ordinary isotopy shows that the definable and PL
images in $\pi_0P_{\mathrm{TOP}}(B)$ coincide.  The same
component obstruction therefore applies.
\end{proof}


\begin{corollary}\label{cor:cyclic-three-absolute-obstruction}
The homeomorphism $F$ in
Theorem~\ref{thm:cyclic-three-concordance-obstruction} can be chosen so
that its component in the absolute ordinary group
$\operatorname{Homeo}(B\times I)$ contains no PL homeomorphism and
no definable homeomorphism, for every fixed o-minimal expansion of
$\mathbb R$.
\end{corollary}

\begin{proof}
Set $X=B\times I$ and fix a compatible smoothing. The homeomorphism
constructed above has nonzero formal Casson--Sullivan class, so its
isotopy class relative to $\partial X$ has no diffeomorphism
representative. It preserves orientation: it fixes each boundary
point and hence has positive orientation sign in a boundary chart.
Galvin--Ladu's boundary-forgetting lemma
\cite[Lemma~5.3, pp.~17--18]{GalvinLadu2025Boundary}
implies that its absolute ordinary isotopy class has no
diffeomorphism representative either. The lemma applies to compact
oriented smooth four-manifolds with boundary, including the two
boundary components here. Its proof uses the boundary restriction
fibrations, injectivity of the three-dimensional boundary comparison
on components, and surjectivity on fundamental groups: the first
makes an absolute smooth representative relative to the boundary,
and the second removes the remaining smooth boundary twist.

Let $H$ be a PL self-homeomorphism of $X$, and write $\sigma$ for
the fixed compatible smoothing. Both $\sigma$ and $H^*\sigma$
are compatible smoothings of the same PL manifold. They are isotopic
as smoothings, by
\cite[Theorem~3.5(2), pp.~17--18; boundary convention p.~8]
{FriedlNagelOrsonPowell2025}.
Here isotopy of smoothings means that there is a diffeomorphism
\[
 u:(X,\sigma)\longrightarrow(X,H^*\sigma)
\]
whose underlying homeomorphism is isotopic to the identity
\cite[Preface, pp.~v--vi]{HirschMazur1974}.
By definition of pullback smoothing, $Hu$ is a diffeomorphism of
$(X,\sigma)$, and an isotopy from the identity to $u$ gives an
ordinary isotopy from $H$ to $Hu$. Thus a PL representative in the
absolute component of $F$ would contradict the preceding conclusion.

Finally, Shiota's given-map isotopy theorem
\cite[Complement to Theorem~2.1, pp.~17--18]{Shiota2010}
isotopes every definable homeomorphism of the compact polyhedron
$X$ to a PL homeomorphism. Over the real field this is also an
ordinary isotopy. Hence the same absolute component contains no
definable representative.
\end{proof}

\paragraph{The finite projective-space example.}
One explicit PL model of $\mathbb {RP}^{3}$ is the quotient
of the barycentric subdivision of the boundary of the
four-dimensional cross-polytope by the antipodal action.
Equivalently, take the order complex of the nonzero partial
sign vectors in $\{-1,0,1\}^4$, modulo simultaneous sign
reversal; order two orbits by compatible strict support
inclusion.  This is a PL quotient of $S^3$ by a free involution,
so it is the standard oriented PL projective three-space and
its fundamental group is $\mathbb Z/2$.

Its face vector is $(40,232,384,192)$.  More precisely, if
$S(r,k)$ denotes a Stirling number of the second kind, then
\[
 f_{k-1}=\sum_{r=k}^{4}\binom4r2^{r-1}k!S(r,k),
 \qquad 1\leq k\leq4.
\]
Choose the top support, its signs up to global reversal, and
the ordered nonempty blocks added along a chain to obtain
this formula.  The usual one-cell-per-dimension CW model has
integral cellular complex
\[
 0\longrightarrow\mathbb Z\xrightarrow{\,0\,}\mathbb Z
 \xrightarrow{\,2\,}\mathbb Z\xrightarrow{\,0\,}\mathbb Z
 \longrightarrow0.
\]
Thus $H_1(\mathbb {RP}^3;\mathbb Z)=\mathbb Z/2$ and
$H^2(\mathbb {RP}^3;\mathbb Z/2)=\mathbb Z/2$, as required.
The finite triangulation verifies the concrete base; the
existence of the obstructed concordance is the geometric
argument above.

\paragraph{Relation to the literature.}
For finite-cyclic elliptic three-manifolds, Galvin already
obtains a related failure of surjectivity on $\pi_1$ for
block diffeomorphism and block homeomorphism spaces
\cite[Theorem~9.0.1, Corollary~9.0.2, and Section~9.3,
pp.~109--117]{GalvinThesis2024}.
The finite cyclic case here uses his relative
Casson--Sullivan realization and makes explicit the passage
to ordinary free-top concordance components, using the
jointly PL replacement of the top path.
